\section{The minimizing-geodesic nerve}
\label{sec:nerve}

Let $X=\widetilde M$ with its lifted negatively curved metric,
and put $B=M\times M$, $\widetilde B=X\times X$.
The group $G=\Gamma\times\Gamma$ acts on $\widetilde B$ by
\[
 (g,h)\cdot(u,v)=(gu,hv)
\]
and on the label set $\Gamma$ by
$(g,h)\cdot a=gah^{-1}$.
For $a\in\Gamma$, let
\[
 L_a(u,v)=d_X(u,av),\qquad
 m(u,v)=\min_{a\in\Gamma}L_a(u,v)=d_M(x,y).
\]
All minima are attained, and the set of minimizing labels
$S(u,v)=\{a:L_a(u,v)=m(u,v)\}$ is finite.

\subsection{The nerve and a uniform allowance}

For every label $a$, define the closed subset
\[
 D_a=\{(u,v)\in\widetilde B:L_a(u,v)=m(u,v)\}.
\]
The \emph{minimizing-geodesic nerve} $\mathcal N_g$
has vertex set $\Gamma$ and a simplex on each finite
nonempty $S\subset\Gamma$ with
$\bigcap_{a\in S}D_a\ne\varnothing$.
Equivariance of $L_a$ gives a simplicial $G$-action.

For $\delta>0$, also define the open sets
\[
 D_a^\delta=\{(u,v):L_a(u,v)<m(u,v)+\delta\}
\]
and let $\mathcal N_g^\delta$ be their nerve.
Clearly $\mathcal N_g\subseteq\mathcal N_g^\delta$.

\begin{lemma}\label{lem:uniform-allowance}
There exists $\delta_0>0$ such that
\[
 \mathcal N_g^\delta=\mathcal N_g
 \quad\text{for every }0<\delta\leq\delta_0.
\]
Moreover, $\mathcal N_g$ has finitely many $G$-orbits of simplices.
\end{lemma}

\begin{proof}
Choose a compact fundamental region $K\subset X$.
For $(u,v)\in K\times K$ and $L_a(u,v)\leq m(u,v)+1$,
the label $a$ belongs to the finite set
\[
 F=\{a\in\Gamma:d_X(K,aK)\leq\operatorname{diam}(M)+1\}.
\]
Thus every simplex of $\mathcal N_g^\delta$,
$0<\delta\leq1$, is $G$-equivalent to a subset
of the fixed finite set $F$.
For a fixed nonempty subset $S\subset F$, put
\[
 \theta_S(u,v)=\max_{a\in S}\bigl(L_a(u,v)-m(u,v)\bigr).
\]
If $S$ does not span a simplex of $\mathcal N_g$,
the continuous function $\theta_S$ is strictly positive
on the compact set $K\times K$, hence has a positive minimum.
There are only finitely many such $S$. Taking $\delta_0$
below all their positive minima proves the equality
of nerves. The same finite-label argument yields the
finiteness of simplex orbits.
\end{proof}

\begin{remark}
The inequality in the definition of $D_a^\delta$ is strict.
Therefore the value of $\delta_0$ can be reduced so that
the proof works without a boundary equality issue.
The lemma concerns the nerves of the covering sets,
not equality of the sets $D_a$ and $D_a^\delta$ themselves.
\end{remark}

\subsection{The equivariant map}

Choose $0<\delta\leq\delta_0$ and define
\[
 w_a(u,v)=\max\{0,m(u,v)+\delta-L_a(u,v)\},
 \qquad W(u,v)=\sum_{a\in\Gamma}w_a(u,v).
\]
At every $(u,v)$ at least one label minimizes $L_a$,
so $W(u,v)>0$.

\begin{proposition}\label{prop:Phi}
The expression
\begin{equation}\label{eq:Phi}
 \Phi_g(u,v)=\sum_{a\in\Gamma}
 \frac{w_a(u,v)}{W(u,v)}[a]
\end{equation}
defines a continuous $G$-equivariant map
\[
 \Phi_g:\widetilde B\longrightarrow|\mathcal N_g|.
\]
It descends to a continuous map
\[
 \bar\Phi_g:B\longrightarrow Q_g=|\mathcal N_g|/G.
\]
\end{proposition}

\begin{proof}
On any compact subset of $\widetilde B$, the label set
for which $w_a$ can be nonzero is finite, by proper
discontinuity of the deck action. Thus \eqref{eq:Phi}
is a locally finite sum of continuous functions.
Its support at $(u,v)$ is a simplex of
$\mathcal N_g^\delta=\mathcal N_g$, hence its image
belongs to the geometric realization.
The identity
\[
 L_{gah^{-1}}(gu,hv)=L_a(u,v)
\]
and invariance of $m$ establish equivariance.
The quotient statement follows.
\end{proof}

\begin{lemma}\label{lem:no-inversions}
The $G$-action on $\mathcal N_g$ is without inversions:
every element preserving a simplex setwise fixes its vertices.
Consequently $Q_g$ has a finite CW structure with one
open cell for each $G$-orbit of simplices.
\end{lemma}

\begin{proof}
Suppose $(g,h)$ preserves a finite vertex set $S$.
For every $a\in S$, some common positive power $r$
of $(g,h)$ fixes $a$, so $g^r=ah^ra^{-1}$.
The torsion-free hyperbolic group $\Gamma$ has the
unique-root property: $p^r=q^r$ implies $p=q$.
It follows that $g=aha^{-1}$ for every $a\in S$,
so $(g,h)$ fixes every vertex.
The standard quotient-cell construction gives the
last assertion, using Lemma~\ref{lem:uniform-allowance}.
\end{proof}

The quotient $Q_g$ need not be contractible. Its relevance
comes instead from the fundamental cycle represented by
$\bar\Phi_g$ and its image in the exterior augmentation
module under Lemma~\ref{lem:cycle-injective}.

\subsection{Multiplicity and the dimension of the nerve}

If $g$ is regular in the sense of Section 1, every minimizing
label set has cardinality at most $2n+1$. Consequently,
\[
 \dim\mathcal N_g\leq2n.
\]
At a point with precisely $k$ minimizing labels, the
length-difference equations cut out a smooth local stratum
of codimension $k-1$ in $B$, by the implicit function
theorem. The inequalities excluding all remaining labels
are strict in a sufficiently small neighborhood.
Thus the local geometry of each simplex is controlled
by the distance-difference functions, even though the
global nerve is defined only in terms of closed sets.

\begin{proposition}\label{prop:top-simplex}
Suppose $g$ is regular and set $N=2n$.
Every pair in $C_g$ determines an $N$-simplex
of $\mathcal N_g$. Conversely, every $N$-simplex
is realized by one or more isolated endpoint pairs.
There are only finitely many such pairs modulo $G$,
and thus $C_g$ is finite.
\end{proposition}

\begin{proof}
An $N$-simplex has $N+1$ minimizing labels.
Its defining $N$ length-difference equations have
full-rank differential at each realization, so their
solutions are isolated. To see that the set of
realizations modulo $G$ is finite, put endpoints
in compact fundamental regions. By properness,
only finitely many labels can minimize there.
If infinitely many pairs realized top simplices,
a convergent subsequence would have a fixed
label configuration and converge to another
simultaneous minimum. Full-rank transversality
at that limit contradicts accumulation.
Each pair in $C_g$ has exactly $N+1$ labels,
so it determines exactly one top simplex after lifts.
Different pairs can determine the same simplex
orbit; no global bijection with simplex orbits
is asserted.
\end{proof}

\begin{remark}\label{rem:nerve-geometry}
The full simplex $Y$ from Section~\ref{sec:canonical}
is contractible and contains $\mathcal N_g$ equivariantly.
The composite
$B\xrightarrow{\bar\Phi_g}Q_g\to Q$
induces, on the fundamental class, the element
$A_M$ of \eqref{eq:A-comparison}. We do not require
the closed sets $D_a$ to constitute a good cover.
\end{remark}
